% Zdroj: PDF priklady-jaro-2026.pdf, fyzická str. 22. Značka: specializace UI-SU, UI-ZPJ. Zařazeno: S3.
\section{Shlukování pomocí $K$-means}
\szzotazka{jaro 2026}{28}{Shlukování pomocí K-means}

\noindent\textbf{Zadání.}\par
\begin{enumerate}
  \item Popište algoritmus $K$-means. K čemu se používá? Popište podrobně jeho kroky a jak je možné algoritmus inicializovat.
  \item Jakou ztrátovou funkci algoritmus optimalizuje?
  \item Konverguje algoritmus? Zformulujte argument nebo důkaz, proč tomu tak je. Jaké to má praktické důsledky?
\end{enumerate}

\medskip
\noindent\textbf{Náčrt řešení.}\par
\begin{itemize}
  \item $K$-means je algoritmus učení bez učitele (neřízené učení) používaný pro shlukování dat do $K$ skupin podle (geometrické)
    podobnosti jejich rysů.
  \item Algoritmus:
    \begin{enumerate}
      \item Inicializace $K$ centroidů (náhodně nebo pomocí metod jako \texttt{kmeans++}).
      \item Algoritmus: Opakuj, dokud nedojde ke konvergenci (žádná změna přiřazení nebo centroidů):
        \begin{enumerate}
          \item Přiřaď každý datový bod k nejbližšímu centroidu.
          \item Aktualizuj centroidy výpočtem průměru všech bodů v každém shluku.
        \end{enumerate}
    \end{enumerate}
  \item \texttt{Kmeans++} inicializace: začni s náhodným centroidem, poté iterativně vybírej nové centroidy s pravděpodobností úměr\-nou
    čtverci vzdálenosti od nejbližšího již existujícího centroidu.
  \item Ztrátová funkce, kterou algoritmus optimalizuje, je součet čtverců vzdáleností mezi datovými body a jejich přiřazenými
    centroidy.
  \item Algoritmus konverguje, protože každá iterace snižuje hodnotu ztrátové funkce, která je omezená zdola nulou. Zaručuje
    pouze konvergenci k lokálnímu minimu, nikoli nutně ke globálnímu minimu.
  \item Praktické důsledky: Volba počátečních centroidů může ovlivnit konečný výsledek shlukování, a proto může být nutné
    provést více běhů s různými inicializacemi, aby se našlo lepší řešení.
\end{itemize}

\medskip
\noindent\textit{Doplňující vysvětlení (nad rámec oficiálního náčrtu):} Samotné klesání zdola omezené funkce ještě nezaručuje, že se algoritmus po konečně mnoha krocích zastaví. Úplný argument: krok přiřazení (každý bod k~nejbližšímu centroidu) ani krok přepočtu (centroid jako průměr, který minimalizuje součet čtverců ve shluku) ztrátu $J$ nezvýší. Hodnota $J$ po přepočtu je určena jen přiřazením bodů do shluků a~různých přiřazení je konečně mnoho (nejvýše $K^N$). Při každé skutečné změně přiřazení $J$ ostře klesne, takže se žádné přiřazení nemůže zopakovat, a~algoritmus se proto po konečně mnoha iteracích zastaví, v~lokálním minimu $J$.
